\chapter{Discussion and Conclusions}

\section{Summary of the Work}

This thesis develops and tests a Bayesian model for reconstructing the temporal
order of the genomic events in a growing tumour. From a single snapshot of
single-cell sequencing data, the model infers the timing of clone-birth and
copy-number events together with the growth rate of each clone. Four pieces of
information feed into the same latent timeline: the number of somatic mutations on
each branch, the way mutations in a copy-number region split into those that arose
before and after the event, the relative abundance of the sampled clones, and the
topology of the phylogenetic tree. The growth rates are given a hierarchical
random-walk prior on the log scale, which lets related clones share information
while still allowing distinct clones to differ, and the event times are
parametrised by stick-breaking so that the tree ordering constraints are always
satisfied.

A large part of the work went into validating the model, in two stages. The model
was first fitted to data generated by its own process, the well-specified case,
where parameter recovery, simulation-based calibration and a set of sensitivity
checks separate the behaviour of the inference from any question of model fit. It
was then applied to synthetic data from ProCESS \parencite{process}, a spatial
tumour simulator whose biology only partly matches the modelling assumptions, so
that the second stage measures robustness to misspecification.

\section{Main Findings}

The two stages give a consistent picture. Under the well-specified process the
event times are recovered accurately and, just as important, with calibrated
uncertainty: their credible intervals reach close to nominal coverage, and
simulation-based calibration \parencite{talts2018validating} gives rank histograms
compatible with uniformity for every parameter, so the posteriors are neither
over- nor under-confident. The growth rates behave differently. Their point
estimates have a larger relative error, but their intervals still cover the truth,
because the model widens the uncertainty rather than settling on a value the data
cannot determine. This is the identifiability limit of a single snapshot, which
constrains the product of a growth rate and an elapsed time much better than either
factor on its own.

The same conclusions hold on the ProCESS data. Both scenarios are read from a
localized biopsy, the adverse case where a contiguous patch of tissue distorts the
clone proportions, so this is a deliberately hard test rather than an idealised one.
Even then the event pseudotimes are recovered with small absolute error
(Table~\ref{tab:sim-bias}), and the only events that fall outside their intervals are
the root events fixed at the origin of time, a boundary effect rather than a real
failure. The growth rates are recovered well on the four-clone chain and lose
precision, though not calibration, on the deeper six-clone tree, in line with what
the well-specified analysis predicted. Because both trees are read under the same
biopsy, the comparison points to the cause: what limits the per-clone growth is the
depth of the lineage, the single-snapshot floor once more, not the way the tissue is
sampled.

\section{Limitations}

A few limitations should be stated plainly. The first is the single-snapshot
identifiability of growth: from one time point the model can give calibrated
intervals on the per-clone growth rates, but it cannot make them arbitrarily
narrow, and on deep lineages the residual uncertainty is large. This floor is a
property of the data, not of the implementation, and the prior cannot remove it:
along the unidentified direction the likelihood is flat, so the posterior simply
follows the prior. A vague, honest prior therefore leaves the floor in place, whereas
a tighter one would not solve the problem but merely assert an answer, narrow yet
possibly wrong and over-confident, a miscalibration that simulation-based calibration
would expose; genuinely shrinking the growth intervals takes more data, not a sharper
prior. Second, the
inference treats the phylogenetic tree and the assignment of copy-number events to
branches as known inputs; these come from upstream reconstruction, and errors in
them would propagate into the timing estimates. Third, ProCESS is itself a model:
real single-cell data carry noise that is absent here, such as sequencing and
copy-number-calling error, doublets and allele dropout, so the robustness shown in
this chapter is necessary but not sufficient evidence of real-world performance.
Finally, a small fraction of the extreme prior draws produced divergent
transitions; these did not affect the validation, but they suggest that a more
conservative sampler adaptation would be sensible for routine use.

\section{Future Work}

The clearest next step follows from the main limitation. The growth--time degeneracy
arises because the tumour is observed only once, so the natural remedy is a
\emph{longitudinal} design that samples it at more than one time. With serial samples
the ratio of a clone's size between two timepoints measures its growth rate on its
own, independently of the elapsed time, and so separates the product $w\cdot t$ into
its factors. One caveat sharpens the requirement: single-cell composition alone is
invariant to the overall scale, so the timepoints must also carry information on the
absolute tumour burden, such as a total cell count or a circulating-tumour-DNA load,
for the growth rates to be pinned down. Given such data, the per-clone growth that a
single snapshot leaves uncertain would become tightly identified. A second, methodological direction is to strengthen the
calibration check with the data-dependent test quantities of
\textcite{modrak2023simulation}, which catch problems that the parameter ranks
alone miss, and to let the molecular clock vary by branch to absorb rate
heterogeneity. A third is to infer the tree and the copy-number assignments jointly
with the timings, instead of fixing them, so that their uncertainty is carried
through. Finally, replacing exponential growth with a spatial or density-dependent
model would bring the model closer to simulators such as ProCESS, and eventually to
real tissue.

\section{Conclusion}

The results show that the evolutionary timing of a tumour can be reconstructed from
single-cell data with a Bayesian model that is accurate and clear about its own
uncertainty. Event times are recovered reliably and with calibrated intervals;
growth rates are recovered as far as a single snapshot allows, with the residual
uncertainty reported rather than hidden; and the inference is robust to the way the
tissue is sampled. The one quantity a single snapshot cannot fully pin down, the
absolute scale of the per-clone growth rates, also points to the clearest way
forward: a \emph{longitudinal} design, in which the same tumour is sequenced at more
than one time and anchored to an absolute measure of its burden, such as a total cell
count or a circulating-tumour-DNA load, would supply precisely the information the
single snapshot lacks. The ratio of a clone's size between two timepoints measures its
growth directly, independently of the elapsed time, so serial data would turn the
growth rates from quantities the snapshot can only bound into sharply identified ones,
while the timing machinery developed here carries over unchanged. The staged
validation, from the well-specified process through
simulation-based calibration to ProCESS, is part of the contribution, since it
gives a way to trust the inferences before any real data are used. The model is a
solid starting point for application to patient single-cell sequencing, where
calibrated timing of clonal events could help us understand how, and how fast,
tumours evolve.
